\section*{Problem 390. Factorizing $n!$ into distinct integers larger than $n$}

\subsection*{1) ROUND-2 OBJECTIVE}

Path (A) is not currently realistic: the limit question
\[
f(n)-2n\sim c\,\frac{n}{\log n}
\]
is still open.

So the correct Round-2 goal is the fail-safe route: extract a rigorous theorem that is strictly stronger than the Round-1 black-box summary.

I will prove an explicit lower bound
\[
f(n)-2n\ge \left(\frac{1}{9}+o(1)\right)\frac{n}{\log n}.
\]
This does not settle the constant $c$, but it gives a concrete asymptotic obstruction, not just the qualitative statement ``of order $n/\log n$.''

\subsection*{2) Round-1 FOUNDATION USED}

I use the following Round-1 reformulation.

Let $f(n)$ be the least $m$ for which there exist distinct integers
\[
n<a_1<\cdots<a_t=m
\]
with
\[
n!=a_1a_2\cdots a_t.
\]

Round~1 rewrote the problem using
\[
(n+1)(n+2)\cdots(2n)=n!\,{2n\choose n}.
\]
If $H=f(n)-2n$, and if we compare an optimal factorization with the full interval $(n,2n]$, then there exist sets
\[
B\subseteq \{n+1,\dots,2n\},\qquad
C\subseteq \{2n+1,\dots,2n+H\}
\]
such that
\[
\prod_{b\in B} b={2n\choose n}\prod_{c\in C} c.
\]
Here $B$ is the set of numbers omitted from $(n,2n]$, and $C$ is the set of numbers inserted beyond $2n$.

Round~1 also recorded the known theorem of Erd\H{o}s, Guy and Selfridge that
\[
f(n)-2n
\]
is of order $n/\log n$.

\subsection*{3) NEW INSIGHT / TOOL (ROUND-2)}

The new tool is a prime-counting argument on the interval
\[
\left(\frac{n}{2},\frac{2n}{3}\right].
\]

For each prime $p$ in this interval, the central binomial coefficient ${2n\choose n}$ contributes one factor of $p$, but inside $(n,2n]$ the only multiples of $p$ are $2p$ and $3p$.
Therefore every such prime forces at least one omitted number carrying an extra factor $2$ or $3$.
This converts a \emph{count of medium primes} into a lower bound on the total $2$-adic and $3$-adic mass that must be supplied by the overhang interval $(2n,2n+H]$.

That mechanism yields the explicit constant $1/9$.

\subsection*{4) ATTACK PLAN (ROUND-2)}

The exact Round-1 gap was that the known lower bound was only quoted abstractly.
I now prove a concrete lower bound by the following chain.

\begin{enumerate}
\item Show that every prime $p\in(n/2,2n/3]$ satisfies
\[
v_p\!\left({2n\choose n}\right)=1.
\]

\item Deduce that for each such prime, at least one of $2p$ or $3p$ must lie in $B$.

\item Count how much $2$-adic and $3$-adic valuation the set $C\subseteq(2n,2n+H]$ can contribute in total.

\item Compare both sides and apply the prime number theorem to the interval $(n/2,2n/3]$.
\end{enumerate}

This gives a fully rigorous explicit asymptotic lower bound for $H=f(n)-2n$.

\subsection*{5) WORK (ROUND-2)}

Let
\[
H=f(n)-2n.
\]
Choose an optimal representation of $n!$ and the corresponding sets $B$ and $C$ from the Round-1 reformulation:
\[
\prod_{b\in B} b={2n\choose n}\prod_{c\in C} c,
\]
with
\[
B\subseteq \{n+1,\dots,2n\},\qquad
C\subseteq \{2n+1,\dots,2n+H\}.
\]

\medskip
\noindent\textbf{Claim 1.}
\emph{For every sufficiently large $n$ and every prime $p$ with}
\[
\frac{n}{2}<p\le \frac{2n}{3},
\]
\emph{one has}
\[
v_p\!\left({2n\choose n}\right)=1.
\]

\medskip
\noindent\textbf{Proof.}
For $n\ge 9$ and $p>n/2$, we have $p^2>2n$.
Hence, by Legendre's formula,
\[
v_p\!\left({2n\choose n}\right)
=
\left\lfloor\frac{2n}{p}\right\rfloor-2\left\lfloor\frac{n}{p}\right\rfloor.
\]
Now
\[
\left\lfloor\frac{n}{p}\right\rfloor=1
\qquad\text{because}\qquad
\frac{n}{2}<p\le n,
\]
and
\[
\left\lfloor\frac{2n}{p}\right\rfloor=3
\qquad\text{because}\qquad
3\le \frac{2n}{p}<4
\]
when $n/2<p\le 2n/3$.
Therefore
\[
v_p\!\left({2n\choose n}\right)=3-2=1.
\]
\hfill\textit{QED.}

\medskip
\noindent\textbf{Claim 2.}
\emph{For every sufficiently large $n$ and every prime $p$ with $n/2<p\le 2n/3$, at least one of $2p$ or $3p$ belongs to $B$.}

\medskip
\noindent\textbf{Proof.}
By Claim~1, the right-hand side
\[
{2n\choose n}\prod_{c\in C} c
\]
is divisible by $p$.

Hence the left-hand side $\prod_{b\in B}b$ is divisible by $p$, so some element of $B$ is divisible by $p$.
But any multiple of $p$ lying in $(n,2n]$ must equal $2p$ or $3p$, because a multiplier $r$ with $rp\in(n,2n]$ satisfies
\[
1<r<4
\]
when $p>n/2$.
Thus at least one of $2p$ or $3p$ lies in $B$.
\hfill\textit{QED.}

\medskip
\noindent\textbf{Claim 3.}
\emph{Let}
\[
P_n:=\left\{p\ \text{prime}:\ \frac{n}{2}<p\le \frac{2n}{3}\right\}.
\]
\emph{Then}
\[
|P_n|\le v_2\!\left(\prod_{b\in B}b\right)+v_3\!\left(\prod_{b\in B}b\right).
\]

\medskip
\noindent\textbf{Proof.}
For each $p\in P_n$, Claim~2 lets us choose a number
\[
b_p\in B\cap\{2p,3p\}.
\]
Each such $b_p$ contributes exactly $1$ to
\[
v_2(b_p)+v_3(b_p),
\]
because $2p$ has $(v_2,v_3)=(1,0)$ and $3p$ has $(v_2,v_3)=(0,1)$.

Also the numbers $b_p$ are distinct.
Indeed, if $b_p=b_q$ for distinct $p,q\in P_n$, then that common number in $(n,2n]$ would be divisible by two distinct primes $>n/2$, impossible because their product would exceed $2n$.

Therefore
\[
|P_n|
=
\sum_{p\in P_n} 1
=
\sum_{p\in P_n}\bigl(v_2(b_p)+v_3(b_p)\bigr)
\le
v_2\!\left(\prod_{b\in B}b\right)+v_3\!\left(\prod_{b\in B}b\right).
\]
\hfill\textit{QED.}

\medskip
\noindent\textbf{Claim 4.}
\emph{One has}
\[
v_2\!\left(\prod_{b\in B}b\right)+v_3\!\left(\prod_{b\in B}b\right)
\le
\frac{3}{2}H+O(\log n).
\]

\medskip
\noindent\textbf{Proof.}
From
\[
\prod_{b\in B} b={2n\choose n}\prod_{c\in C} c
\]
we get
\[
v_2\!\left(\prod_{b\in B}b\right)+v_3\!\left(\prod_{b\in B}b\right)
=
v_2\!\left({2n\choose n}\right)+v_3\!\left({2n\choose n}\right)
+
v_2\!\left(\prod_{c\in C}c\right)+v_3\!\left(\prod_{c\in C}c\right).
\]

First, the central binomial coefficient contributes only $O(\log n)$ to the $2$- and $3$-adic valuations.
Indeed, by Legendre's formula,
\[
v_p\!\left({2n\choose n}\right)
=
\sum_{j\ge 1}\left(\left\lfloor\frac{2n}{p^j}\right\rfloor-2\left\lfloor\frac{n}{p^j}\right\rfloor\right),
\]
and each summand is at most $1$, while all summands vanish once $p^j>2n$.
Hence
\[
v_2\!\left({2n\choose n}\right)\le \lfloor \log_2(2n)\rfloor+1,
\qquad
v_3\!\left({2n\choose n}\right)\le \lfloor \log_3(2n)\rfloor+1,
\]
so
\[
v_2\!\left({2n\choose n}\right)+v_3\!\left({2n\choose n}\right)=O(\log n).
\]

Second, because $C\subseteq \{2n+1,\dots,2n+H\}$,
\[
v_2\!\left(\prod_{c\in C}c\right)
\le
\sum_{m=2n+1}^{2n+H} v_2(m)
=
\sum_{j\ge 1} \#\{m\in[2n+1,2n+H]:2^j\mid m\}.
\]
For each $j$, this count is at most $H/2^j+1$, and only $O(\log n)$ values of $j$ can occur because $2^j\le 2n+H$ is necessary for the count to be nonzero.
Hence
\[
v_2\!\left(\prod_{c\in C}c\right)
\le
\sum_{j\ge 1}\left(\frac{H}{2^j}+1\right)
=
H+O(\log n).
\]
Similarly,
\[
v_3\!\left(\prod_{c\in C}c\right)
\le
\sum_{j\ge 1}\left(\frac{H}{3^j}+1\right)
=
\frac{H}{2}+O(\log n),
\]
again because only $O(\log n)$ values of $j$ contribute.

Adding these bounds gives
\[
v_2\!\left(\prod_{b\in B}b\right)+v_3\!\left(\prod_{b\in B}b\right)
\le
\frac{3}{2}H+O(\log n).
\]
\hfill\textit{QED.}

\medskip
\noindent\textbf{Claim 5.}
\emph{For sufficiently large $n$,}
\[
H\ge \left(\frac{1}{9}+o(1)\right)\frac{n}{\log n}.
\]

\medskip
\noindent\textbf{Proof.}
By Claim~3 and Claim~4,
\[
|P_n|\le \frac{3}{2}H+O(\log n).
\]
By the prime number theorem,
\[
|P_n|
=
\pi\!\left(\frac{2n}{3}\right)-\pi\!\left(\frac{n}{2}\right)
=
\left(\frac{2}{3}-\frac{1}{2}+o(1)\right)\frac{n}{\log n}
=
\left(\frac{1}{6}+o(1)\right)\frac{n}{\log n}.
\]
Hence
\[
\left(\frac{1}{6}+o(1)\right)\frac{n}{\log n}
\le
\frac{3}{2}H+O(\log n).
\]
Since $O(\log n)=o(n/\log n)$, division by $3/2$ gives
\[
H\ge \left(\frac{1}{9}+o(1)\right)\frac{n}{\log n}.
\]
Because $H=f(n)-2n$, this is exactly
\[
f(n)-2n\ge \left(\frac{1}{9}+o(1)\right)\frac{n}{\log n}.
\]
\hfill\textit{QED.}

\subsection*{6) ADVERSARIAL VERIFICATION}

I check the vulnerable points.

\begin{itemize}
\item \textbf{Did I use the Round-1 set identity correctly?}
Yes.  The identity
\[
\prod_{b\in B} b={2n\choose n}\prod_{c\in C} c
\]
is exactly the Round-1 reformulation, and every valuation argument starts from it.

\item \textbf{Could a prime $p\in(n/2,2n/3]$ be absorbed by some number in $(n,2n]$ other than $2p$ or $3p$?}
No.  Multipliers $1$ and $4$ are impossible because $p\le n$ and $4p>2n$ in this interval.  So only $2p$ and $3p$ occur.

\item \textbf{Could two different primes share the same omitted number?}
No.  Any number in $(n,2n]$ has at most one prime factor greater than $n/2$, because two such primes would have product $>2n$.

\item \textbf{Does $C$ perhaps contribute more $2$- and $3$-adic mass than estimated?}
No.  I bounded the valuation of \emph{the entire interval} $[2n+1,2n+H]$, so the bound automatically applies to its subset $C$.

\item \textbf{Small $n$.}
Claim~1 was stated only for sufficiently large $n$ because I used $p^2>2n$.
That is harmless: the final statement is asymptotic, and finitely many small values do not affect the $o(1)$ term.

\item \textbf{Interaction with the known upper bound.}
Round~1 already cited the Erd\H{o}s--Guy--Selfridge upper bound
\[
f(n)-2n=O\!\left(\frac{n}{\log n}\right).
\]
The new theorem is compatible with that and strengthens only the lower-bound side.
\end{itemize}

\subsection*{7) FINAL}

\textbf{UNRESOLVED (BUT STRICTLY ADVANCED).}

The strict Round-2 advance is the explicit theorem
\[
f(n)-2n\ge \left(\frac{1}{9}+o(1)\right)\frac{n}{\log n}.
\]

Thus the previously quoted qualitative lower bound has been replaced by a concrete asymptotic constant.
Combined with the known upper bound of Erd\H{o}s--Guy--Selfridge,
\[
f(n)-2n=O\!\left(\frac{n}{\log n}\right),
\]
we now have a sharper quantitative bracket.

What remains open is the actual question of the problem:
does the normalized quantity
\[
\frac{(f(n)-2n)\log n}{n}
\]
converge, and if so to what constant $c$?

\subsection*{8) COMPLETION ESTIMATE (MANDATORY)}

COMPLETION: 60\%

\subsection*{9) REFERENCES}

\begin{enumerate}
\item P. Erd\H{o}s, R. K. Guy, and J. L. Selfridge, \emph{The product of the divisors of an integer}, for the classical upper and lower bounds showing that $f(n)-2n$ is of order $n/\log n$.

\item The Erd\H{o}s Problems site, Problem~390, for the statement of the problem and the citation to the Erd\H{o}s--Guy--Selfridge result.

\item The prime number theorem, used in the standard form
\[
\pi(an)-\pi(bn)=\left((a-b)+o(1)\right)\frac{n}{\log n}
\]
for fixed $0<b<a$.
\end{enumerate}
